\documentclass[12pt]{article}

\input{../../_assets/preambulo.tex}


\begin{document}

    % 1. Foto de fondo
    % 2. Título
    % 3. Encabezado Izquierdo
    % 4. Color de fondo
    % 5. Coord x del titulo
    % 6. Coord y del titulo
    % 7. Fecha

    
    \input{../../_assets/portada}
    \portadaExamen{ffccA4.jpg}{Ecuaciones\\Diferenciales II\\Examen XIX}{Ecuaciones Diferenciales II. Examen XIX}{MidnightBlue}{-8}{28}{2026}{}

    \begin{description}
        \item[Asignatura] Ecuaciones Diferenciales II.
        \item[Curso Académico] 2014/2015.
        \item[Grado] Doble Grado en Ingeniería Informática y Matemáticas.
        \item[Grupo] Único.
        \item[Profesor] Rafael Ortega Ríos.
        \item[Descripción] Convocatoria Extraordinaria.
        \item[Fecha] 17 de septiembre de 2015.
        % \item[Duración] ---.
    
    \end{description}
    \newpage


    % ------------------------------------
    
    \begin{ejercicio}
        Calcula la solución maximal del problema
        \[ \dot{x} = x^2 + 1, \quad x(0) = 1. \]
    \end{ejercicio}

    \begin{ejercicio}
        Encuentra una retracción del plano en el conjunto
        \[ E = \left\{ (x,y) \in \bb{R}^2 : x^2 + 4y^2 \leq 1 \right\}. \]
    \end{ejercicio}

    \begin{ejercicio}
        Encuentra una sucesión de funciones continuas $f_n : \left[0, 1\right] \to \bb{R}$ de manera que $\{f_n\}$ sea uniformemente acotada pero no sea equicontinua.
    \end{ejercicio}

    \begin{ejercicio}
        Se considera la sucesión de funciones $\{x_n\}$, $x_n : \left[0, 1\right] \to \bb{R}$ definidas de manera inductiva por las fórmulas
        \[ x_0(t) = 1, \quad x_{n+1}(t) = 1 + \int_0^t \sen(x_n(\tau)) \, d\tau. \]
        Demuestra que existe una sucesión parcial que converge uniformemente.
    \end{ejercicio}

    \begin{ejercicio}
        Dado $n \geq 1$ se considera la solución $x_n(t)$ del problema
        \[ \dot{x} = \cos\left(\nicefrac{x}{n}\right), \quad x(0) = 1. \]
        Calcula, si existe, $\lim\limits_{n\to\infty} x_n(t)$.
    \end{ejercicio}

    \begin{ejercicio}
        Se considera la ecuación $\ddot{x} + \dot{x} + \sen x = 0$. ¿Es el equilibrio $x = 0$ asintóticamente estable?
    \end{ejercicio}

    \begin{ejercicio}
        Encuentra un plano que sea invariante para el sistema $\dot{x} = Ax$ donde $A$ es la matriz
        \[ \begin{pmatrix} 3 & 0 & 0 \\ 0 & 3 & 1 \\ 0 & -1 & 3 \end{pmatrix}. \]
    \end{ejercicio}

    \begin{ejercicio}
        Se supone que $f : \bb{R}^2 \to \bb{R}$, $(t,x) \mapsto f(t,x)$ es continua y cumple
        \[ |f(t,x) - f(t,y)| \leq 7|x - y|. \]
        Demuestra que la ecuación integral
        \[ x(t) = 2 + \int_0^t \frac{1}{\sqrt{|s|}} f(s, x(s)) \, ds \]
        tiene a lo sumo una solución continua y definida en todo $\bb{R}$.
    \end{ejercicio}

    \newpage
    \setcounter{ejercicio}{0} % Reiniciar contador de ejercicios
    % \noindent
    % \textbf{Solución.}

\end{document}